% =====================================================================
%  CHAPTER 7: 行为心理因素与健康（对应大纲第七章）
%  参考PPT：第七章 行为心理因素与健康
% =====================================================================
\chapter{行为心理因素与健康}
\setcounter{chapter}{7}
\chapterintro{
掌握行为及健康相关行为的概念；掌握人格与健康（A/B/C型性格）；熟悉现代压力基本理论（生物应激/社会事件刺激/心理认知/现代压力理论）和持久过大压力的不良后果；熟悉健康相关行为的基本特征；熟悉吸烟的危害与控烟政策措施（FCTC、MPOWER）；了解行为心理问题的干预方法（政策/环境/媒体/社区/组织干预）。
}

\section{心理因素与健康}

\subsection{人格与健康}

\begin{defbox}
\textbf{人格（Personality）：}是稳定地表现于个体的心理特质，由遗传和环境共同决定，对人的社会化过程的作用举足轻重。

\textbf{人格与健康：}消极人格（抑郁、愤怒、敌意与焦虑）→较低健康水平；刚性人格（责任感、控制感、挑战性）→较好健康水平。

\textbf{性格类型与疾病：}
\begin{itemize}[leftmargin=*]
    \item \imp{A型性格}：急躁、缺乏耐性，具有敌意、匆忙和竞争的特点，与冠心病、高血压、溃疡病等有关
    \item \imp{B型性格}：不争强好胜、做事不慌不忙
    \item \imp{C型性格}：压抑自己的情绪，过分忍让，回避矛盾，怒而不发，好生闷气，内向——肿瘤的性格模型
\end{itemize}
\end{defbox}

\subsection{心理压力与健康}

\begin{keybox}
\textbf{压力（Stress）/ 心理应激：}是指人们生活中的各种刺激事件（应激源）和内在要求在心理上所构成的困惑或威胁，表现为心身紧张或不适。

\textbf{压力基本理论：}
\begin{enumerate}[leftmargin=*]
    \item \imp{生物应激理论}：应激反应是生物体对某些情景自动的或激素反应（克劳德：机体对外界刺激所做出的适应性反应）
    \item \imp{社会事件刺激理论}：压力是外部压力事件的刺激作用（坎农）
    \item \imp{心理认知理论}：拉扎勒斯和福尔克曼强调人们对刺激事件的认知和解释的重要性
    \item \imp{现代压力理论（应激过程模型）}：将心理应激看作是由各种刺激事件（应激源）到应激反应的多因素作用过程
\end{enumerate}

\textbf{应激源分类：}(1)躯体性应激（直接对机体作用）；(2)心理性应激源（人际冲突、个体强烈需求、能力不足与期望矛盾）；(3)社会性应激源（客观社会学指标、社会变动与地位不合适）；(4)文化性应激源（语言/风俗/习惯改变）。

\textbf{持久过大压力的不良后果：}(1)健康问题——各类疾病（高血压、偏头疼、癌症、溃疡、关节炎、心血管病等），心理和行为问题（心理障碍、吸烟、酗酒、自杀和反社会行为）；(2)工作问题——旷工、消极怠工、工作差错和事故增加；(3)管理和决策问题——决策失误（建设性思维减少/扭曲性思维增加、危险性抉择增加、重视短期忽略长期目标）。
\end{keybox}

\section{行为生活方式与健康}

\begin{keybox}
\textbf{行为（Behavior）：}人类在内外因素共同作用下产生的外部活动。本能行为（先天性的，摄食/睡眠/性等）+ 社会行为（在社会文化环境中通过社会化过程获得的，锻炼/吸烟/酗酒等）。

\textbf{健康相关行为（Health Related Behavior）：}指个体或群体与健康和疾病有关的行为。分为：
\begin{itemize}[leftmargin=*]
    \item \imp{促进健康的行为（Health-Promoted Behavior）}：个体或群体表现出的、客观上有利于自身和他人健康的行为
    \item \imp{危害健康的行为（Health-Risky Behavior）}：偏离个人、他人和社会健康期望、不利于健康的行为
\end{itemize}

\textbf{促进健康行为的基本特征：}有利性、规律性、和谐性、一致性、适宜性。

\textbf{行为因素与健康的关系：}2008年WHO调查显示，50\%的死亡是由于行为生活方式因素、30\%为环境因素、10\%为生物遗传因素、10\%为医疗卫生服务因素。发达国家的经验证实，通过行为干预，可使高血压发病率下降55\%，脑卒中下降75\%，糖尿病下降50\%，肿瘤下降33\%。
\end{keybox}

\section{行为心理问题的干预}

\begin{defbox}
\textbf{（一）政策干预}——效益很高的措施，如公共场所禁烟计划、禁止吸烟广告。政府是政策制定的主体，需通过政策倡导促动来促使政府采纳和实施有利于人民健康的政策。\textbf{倡导促动（PCPA模式）：}倡议（提出议题）→联盟（建立伙伴关系）→宣传（公开意图、设计讯息）→行动（游说和鼓动）。

\textbf{（二）环境工程设施干预}——通过改善环境工程设施增加行为相关物的可得性或可及性，以增强行为干预效果。

\textbf{（三）大众媒体干预}——利用媒体影响大众和决策者的知识、观点、态度和行为，进行健康教育。

\textbf{（四）社区干预}——在社区层面进行疾病预防、控制、教育和咨询等活动。社区动员是社区干预的前提。

\textbf{（五）组织干预}——通过对不合理的组织结构和行为进行改变达到干预目标，分4个阶段：问题知觉→行动启动（采纳）→实施→定型化。
\end{defbox}

\section{常见健康相关行为——吸烟与控烟}

\begin{keybox}
\textbf{吸烟的危害：}心脑血管疾病、呼吸系统疾病、消化系统疾病、生殖系统疾病、神经系统疾病、内分泌系统疾病。我国目前吸烟人数已超过3亿，遭受二手烟危害的非吸烟人群高达7.4亿，每年有120多万人死于烟草相关的疾病。

\textbf{控烟的政策措施：}
\begin{itemize}[leftmargin=*]
    \item \imp{《烟草控制框架公约》（FCTC，2003年）}——WHO首次制定的具有国际法约束力的全球公约。中国2005年批准，2006年生效
    \item \imp{MPOWER六项政策（WHO 2008年）：}(1)监测烟草使用与预防政策（Monitor）；(2)保护人们免受烟雾危害（Protect）；(3)提供戒烟帮助（Offer）；(4)警示烟草危害（Warn）；(5)禁止烟草广告、促销和赞助（Enforce）；(6)增加税收和提高烟价（Raise）
\end{itemize}

\textbf{控烟倡导促动：}政策倡导促动（针对决策者的社会行动，促使其出台控烟政策）+ 大众倡导促动（大众宣传+大众劝说）。
\end{keybox}

\section{复习题}

\begin{enumerate}[leftmargin=*, itemsep=8pt]
    \item \textbf{【名词解释】}人格；心理压力（心理应激）；健康相关行为；倡导促动（PCPA模式）
    \item \textbf{【简答题】}简述A/B/C型性格的特征及其与疾病的关系。
    \item \textbf{【简答题】}试述现代压力理论（应激过程模型）及四种应激源。
    \item \textbf{【简答题】}简述促进健康行为的基本特征。
    \item \textbf{【简答题】}简述WHO MPOWER控烟六项政策措施。
\end{enumerate}

\subsection*{参考答案要点}

\begin{enumerate}[leftmargin=*, itemsep=5pt]
    \item \textbf{人格}：稳定表现于个体的心理特质，由遗传和环境共同决定。\textbf{心理压力/应激}：生活中的各种刺激事件和内在要求在心理上构成的困惑或威胁。\textbf{健康相关行为}：个体或群体与健康和疾病有关的行为，分促进健康行为和危害健康行为。\textbf{PCPA模式}：倡议—联盟—宣传—行动的倡导促动模式。
    \item A型：急躁/缺乏耐性/敌意/匆忙/竞争→冠心病/高血压/溃疡病。B型：不争强好胜/做事不慌不忙。C型：压抑情绪/过分忍让/回避矛盾/怒而不发/内向→肿瘤的性格模型。
    \item 应激过程模型：应激源→认知评价（受个性特征影响）→应对方式+社会支持→应激反应→适应或适应不良（身心疾病）。四种应激源：躯体性、心理性（人际冲突/需求矛盾）、社会性（社会学指标变化/社会地位变化）、文化性（语言/风俗/习惯改变）。
    \item 五大特征：(1)有利性（有益于自己/他人/社会）；(2)规律性（有恒常规律）；(3)和谐性（有鲜明个性又能根据环境调整）；(4)一致性（外显行为与内心一致）；(5)适宜性（行为强度有理性控制，无明显冲动）。
    \item MPOWER六项：(1)监测烟草使用与预防政策(Monitor)；(2)保护人们免受烟雾危害(Protect)；(3)提供戒烟帮助(Offer)；(4)警示烟草危害(Warn)；(5)禁止烟草广告/促销/赞助(Enforce)；(6)增加税收和提高烟价(Raise)。
\end{enumerate}

\newpage
